Sokoban is a widely successful puzzle game where the goal is to push boxes around in a warehouse to position the boxes in correct places. However, the positions of the boxes determine the mobility of other boxes and the player. In some cases, the player may accidentally push a box into a corner, rendering it completely immobile \cite{junghanns1999pushing}. Due to these restrictions, players are required to plan ahead of moving to successfully solve each level.

The game was invented in 1981 by Hiroyuki Imabayashi as a hobby and was commercially released in 1982. Since its release, the series has sold over 4.1 million copies worldwide, with over 40 official games and countless clones, enabled by the simplistic gameplay design.

\section{Gameplay}
Sokoban is a top-down puzzle game in which the player controls a 2D character in a warehouse with several boxes that the player can push. Some tiles are marked as storage locations where the boxes in the level need to be placed. In order to complete the level, all boxes need to be placed on designated tiles \cite{murase1996automatic}. The levels are generally surrounded by tight walls, and the boxes are as large as the player. This often means that the boxes can block paths, and the player often needs to relocate a box in order to reach certain areas in the level. This is made more difficult by the fact that the player can only push the boxes and cannot pull them. If a box is pushed against a wall, the player can only move the box tangent to the wall, as they cannot push away from the wall \cite{sokoban_official_whatis}. If the box is pushed into a corner, the box cannot be moved for the same reason \cite{micomgames1983sokoban}. If this is not part of the solution, the player reaches a dead end, and the only way to progress is to reset the current level or undo the previous move.

\section{History}
The first iteration of the game ran on the NEC-8001 as a hobby project by Hiroyuki Imabayashi \cite{login1983thinkingrabbit}. The game was primarily inspired by \textit{Aldebaran \#1}, created by Hudson Soft, in which the player pushed objects to block radiation \cite{monthlymycom1980aldebaran, aguilera2020imabayashi}. However, Imabayashi envisioned a more difficult game in which objects could become stuck at specific places, a game mechanic less explored in \textit{Aldebaran \#1} due to its spacious levels \cite{sokoban_official_greetings}. The hobby project became a commercial project when a salesman came across it in a store owned by Imabayashi's parents-in-law. Since its original launch on the NEC PC-8801 \cite{popcom1984visiting}, the game has been re-released many times across various platforms, from handheld gaming consoles like the Game Boy \cite{atelier_double_history_1989} to consumer computers like the Apple II \cite{ellison1988whyjapan} and Windows \cite{sokoban_smart_official}, and most recently the Nintendo Switch and PlayStation 4 \cite{sokoban_official_lineup}.

\section{Reception}
Sokoban instantly became the favorite of puzzle enthusiasts \cite{tanaka1986thesokoban}, with over 25,000 copies sold by July 1984 following the initial release in Japan \cite{popcom1984strategy}. It was quickly followed by multiple successful releases in Japan, which prompted a global release. The U.S. release \textit{Soko-Ban} sold over 50,000 copies by mid-September 1988 \cite{ellison1988whyjapan}. By 2018, the series had sold over 4.1 million copies worldwide \cite{atpress2018chukyo}.

Sokoban was praised for its compelling, simple gameplay combined with difficult puzzles requiring a high level of planning, comparable in depth to Shogi \cite{micomgames1983sokoban}. The popularity and difficulty combined soon led to a study on understanding the nature of Sokoban as a problem from the computational complexity theory point of view.

The difficulty was soon studied with computational complexity theory by Fryers et al., and the result was published in 1995 in \textit{Eureka} \cite{fryers1995sokoban}. Sokoban was found to be NP-hard, which indicates that any NP problem can be encoded into a Sokoban puzzle in polynomial time. In other words, theoretically, finding a generic solution to all Sokoban puzzles indicates that the logic required to solve the puzzle can be applied to any other easier NP problem, proving that any NP problem is also a P problem, which is tractable to solve. While it is widely believed that NP is not P, searching for a solution to an NP-hard problem has the academic significance that its solution may be applied to some other specific problems. Similarly to how AlexNet triggered the Deep Learning revolution in the 2010s while attempting to solve an image classification problem \cite{krizhevsky2012imagenet}, academic inquiry continually seeks compelling problems whose solutions yield insights that extend beyond their original scope \cite{lecun2015deep}. For this reason, Sokoban is an academically interesting problem. The next chapter discusses the variety of methodologies that researchers have been applying to solve Sokoban.


